\documentclass[12pt]{article}
\usepackage[utf8]{inputenc}
\usepackage[T1]{fontenc}
\usepackage[margin=1in]{geometry}
\usepackage{setspace}
\usepackage{xcolor}
\usepackage[colorlinks=true, urlcolor=blue!60!black, linkcolor=blue!60!black]{hyperref}
\setlength{\parindent}{2em}
\setlength{\parskip}{1.25em}
\renewcommand{\baselinestretch}{1.0}
\begin{document}
\begin{center}
{\LARGE \textit{Kennedy v. Bremerton School District}, 597 U.S. 507 (2022)\par}
\vspace{0.8em}
{\large Case Brief\par}
\end{center}
\vspace{1em}

\section{Memory Jogger}
A public high school football coach's brief, quiet postgame prayer at midfield is protected by the Free Exercise and Free Speech Clauses, and the Establishment Clause---interpreted by historical practices and understandings rather than the abandoned \textit{Lemon} and endorsement tests---neither requires nor permits a school district to suppress it (\href{https://www.courtlistener.com/opinion/6619004/kennedy-v-bremerton-school-dist/}{CourtListener}; \href{https://www.law.cornell.edu/supremecourt/text/21-418}{Cornell LII}).

\section{Detailed Case Facts}
Joseph A. Kennedy, a practicing Christian and Marine Corps veteran, began working in 2008 as a part-time assistant varsity football coach and head junior varsity football coach at Bremerton High School, a public school in Bremerton, Washington (\href{https://www.law.cornell.edu/supremecourt/text/21-418}{Cornell LII}). From the start of his tenure, Kennedy followed a personal practice of kneeling at the 50-yard line after players and coaches shook hands at the end of each game to offer a quiet prayer of thanksgiving lasting approximately 30 seconds (\href{https://www.law.cornell.edu/supremecourt/text/21-418}{Cornell LII}). Initially, Kennedy prayed alone, but over time players asked to join him; he responded, ``This is a free country. You can do what you want,'' and eventually most of the team, and sometimes opposing players, joined him (\href{https://www.law.cornell.edu/supremecourt/text/21-418}{Cornell LII}). Kennedy's postgame practice evolved into short motivational speeches containing overtly religious references while players knelt around him, and he also participated in a pre-existing school tradition of locker-room prayers (\href{https://www.law.cornell.edu/supremecourt/text/21-418}{Cornell LII}).

In September 2015, after an opposing school's employee complimented the practice to Bremerton's principal, the Bremerton School District (the District) initiated an inquiry (\href{https://www.law.cornell.edu/supremecourt/text/21-418}{Cornell LII}). On September 17, 2015, the superintendent sent Kennedy a letter identifying the postgame religious talks and locker-room prayers as problematic under the Establishment Clause and directing that any religious expression while on duty be ``nondemonstrative'' to avoid the perception of school endorsement (\href{https://www.law.cornell.edu/supremecourt/text/21-418}{Cornell LII}). Kennedy complied by ending locker-room prayers and removing religious references from postgame team talks, and after one game he returned to the empty stadium to pray alone (\href{https://www.law.cornell.edu/supremecourt/text/21-418}{Cornell LII}).

On October 14, 2015, Kennedy's attorneys notified the District that his sincerely held religious beliefs compelled him to offer a short, private, personal postgame prayer at midfield and announced he would resume doing so after the October 16 homecoming game (\href{https://www.law.cornell.edu/supremecourt/text/21-418}{Cornell LII}). Kennedy also gave local media interviews publicizing his intent (\href{https://www.law.cornell.edu/supremecourt/text/21-418}{Cornell LII}; \href{https://www.oyez.org/cases/2021/21-418}{Oyez}). On October 16, the District replied with an ultimatum barring ``any overt actions'' appearing to a reasonable observer to endorse prayer while on duty (\href{https://www.law.cornell.edu/supremecourt/text/21-418}{Cornell LII}). After the October 16, 23, and 26 games, while Bremerton players were engaged in other postgame activities such as singing the school fight song, Kennedy knelt at midfield to pray (\href{https://www.law.cornell.edu/supremecourt/text/21-418}{Cornell LII}). No Bremerton players joined him during those three prayers, though opposing players and spectators rushed the field on October 16 (prompting police security at later games), he prayed alone on October 23, and members of the public and state representatives gathered around him on October 26 (\href{https://www.law.cornell.edu/supremecourt/text/21-418}{Cornell LII}). On October 28, 2015, the District placed Kennedy on paid administrative leave, noting in a public Q\&A that it had no evidence students were directly coerced to pray with him, and later issued a negative evaluation advising against rehiring him for failing to follow District policy and failing to supervise student-athletes after games (\href{https://www.law.cornell.edu/supremecourt/text/21-418}{Cornell LII}). Kennedy did not return for the following season.

\section{Procedural History}
Kennedy sued the District in the U.S. District Court for the Western District of Washington, alleging violations of the First Amendment's Free Speech and Free Exercise Clauses (as incorporated against the states via the Fourteenth Amendment) and Title VII of the Civil Rights Act of 1964 (\href{https://www.oyez.org/cases/2021/21-418}{Oyez}; \href{https://www.law.cornell.edu/supremecourt/text/21-418}{Cornell LII}). The district court denied his motion for a preliminary injunction, and the U.S. Court of Appeals for the Ninth Circuit affirmed in \textit{Kennedy v. Bremerton School District}, 869 F.3d 813 (9th Cir. 2017) (\href{https://www.law.cornell.edu/supremecourt/text/21-418}{Cornell LII}). The Supreme Court denied certiorari in 2019, though Justice Alito (joined by Justices Thomas, Gorsuch, and Kavanaugh) issued a statement respecting the denial expressing concern over the Ninth Circuit's reasoning (\href{https://www.law.cornell.edu/supremecourt/text/21-418}{Cornell LII}).

On remand, following discovery, both sides moved for summary judgment. Finding that the District's sole reason for suspending Kennedy was its perceived risk of Establishment Clause liability for the three October 2015 prayers, the district court nonetheless granted summary judgment to the District, 443 F. Supp. 3d 1223 (W.D. Wash. 2020) (\href{https://www.law.cornell.edu/supremecourt/text/21-418}{Cornell LII}). The Ninth Circuit affirmed, 991 F.3d 1004 (9th Cir. 2021), holding that Kennedy's prayers were unprotected government speech and, alternatively, that avoiding an Establishment Clause endorsement violation was a compelling interest justifying his suspension (\href{https://www.law.cornell.edu/supremecourt/text/21-418}{Cornell LII}). The Ninth Circuit denied rehearing en banc over the dissents of 11 judges, 4 F.4th 910 (9th Cir. 2021) (\href{https://www.law.cornell.edu/supremecourt/text/21-418}{Cornell LII}). The Supreme Court granted certiorari on January 14, 2022, heard oral argument on April 25, 2022, and reversed on June 27, 2022 (\href{https://www.oyez.org/cases/2021/21-418}{Oyez}; \href{https://www.supremecourt.gov/opinions/21pdf/21-418_i425.pdf}{Slip Opinion}).

\section{Judicial Votes}
The Supreme Court decided the case by a \textbf{6--3} vote (\href{https://www.law.cornell.edu/supremecourt/text/21-418}{Cornell LII}; \href{https://www.oyez.org/cases/2021/21-418}{Oyez}):
\begin{itemize}
    \item \textbf{Majority (6):} Justice Neil Gorsuch authored the opinion of the Court, joined in full by Chief Justice John Roberts and Justices Clarence Thomas, Samuel Alito, and Amy Coney Barrett, and joined except as to Part III--B (the threshold \textit{Pickering}--\textit{Garcetti} free speech analysis) by Justice Brett Kavanaugh (\href{https://www.law.cornell.edu/supremecourt/text/21-418}{Cornell LII}).
    \item \textbf{Concurrences (2):} Justice Clarence Thomas filed a concurring opinion; Justice Samuel Alito filed a concurring opinion (\href{https://www.law.cornell.edu/supremecourt/text/21-418}{Cornell LII}).
    \item \textbf{Dissent (3):} Justice Sonia Sotomayor filed a dissenting opinion, joined by Justices Stephen Breyer and Elena Kagan (\href{https://www.law.cornell.edu/supremecourt/text/21-418}{Cornell LII}).
\end{itemize}

\section{Holding}
\textbf{Judgment:} Reversed; Kennedy was entitled to summary judgment on his First Amendment Free Exercise and Free Speech claims (\href{https://www.law.cornell.edu/supremecourt/text/21-418}{Cornell LII}).

\textbf{Rule Announced:} The Free Exercise and Free Speech Clauses of the First Amendment protect a public school employee's brief, quiet, personal religious observance undertaken during a postgame lull when employees are free to attend to personal matters, and the Establishment Clause---which must be interpreted by ``reference to historical practices and understandings'' rather than the abandoned \textit{Lemon} and endorsement tests---neither requires nor permits a school district to suppress such private religious expression in the absence of coercion (\href{https://www.law.cornell.edu/supremecourt/text/21-418}{Cornell LII}).

\textbf{Holding vs.\ Dicta:} The Court's core holdings are that (1) the District's non-neutral, non-generally-applicable restriction on Kennedy's three October 2015 prayers burdened his Free Exercise rights; (2) under \textit{Pickering} and \textit{Garcetti} (joined by five Justices in Part III--B), those prayers were private speech rather than government speech; (3) \textit{Lemon v. Kurtzman}, 403 U.S. 602 (1971), and its endorsement offshoot are abandoned in favor of a ``historical practices and understandings'' standard; and (4) the District failed to show coercion or any other Establishment Clause violation under any level of scrutiny (\href{https://www.law.cornell.edu/supremecourt/text/21-418}{Cornell LII}). By contrast, the Court's observations about whether \textit{Employment Division v. Smith} should be reconsidered, whether the Free Exercise Clause might require a different test than \textit{Pickering}--\textit{Garcetti} for public employees, and how a school might treat a Muslim teacher's headscarf or a Christian aide's lunchtime prayer are illustrative dicta or explicitly reserved questions (\href{https://www.law.cornell.edu/supremecourt/text/21-418}{Cornell LII}).

\section{Analysis of Opinions}

\subsection{Majority Opinion (Justice Gorsuch)}
Justice Gorsuch began by framing the Free Exercise and Free Speech Clauses as working ``in tandem'' to ``doubly protect'' expressive religious activity (\href{https://www.law.cornell.edu/supremecourt/text/21-418}{Cornell LII}). Under the Free Exercise Clause, Kennedy easily met his threshold burden: his sincere religious practice of offering a brief, quiet postgame prayer was burdened by District policies that were neither neutral (because prohibiting religious conduct was their admitted object) nor generally applicable (because other coaches were permitted to attend to personal matters like phone calls or greeting friends during the same postgame period) (\href{https://www.law.cornell.edu/supremecourt/text/21-418}{Cornell LII}).

In Part III--B (joined by five Justices), the Court held that Kennedy also carried his threshold burden under the Free Speech Clause's \textit{Pickering}--\textit{Garcetti} framework. Applying \textit{Lane v. Franks}, 573 U.S. 228 (2014), the Court concluded that Kennedy's prayers did not ``owe their existence'' to his coaching duties: he was not instructing players, calling plays, or delivering a government-commissioned message, and he prayed during a lull while students were singing the fight song or otherwise occupied (\href{https://www.law.cornell.edu/supremecourt/text/21-418}{Cornell LII}). Treating everything a coach says while on duty as government speech would posit an ``excessively broad job description'' and allow schools to ban a Muslim teacher's headscarf or a teacher's quiet lunchtime prayer (\href{https://www.law.cornell.edu/supremecourt/text/21-418}{Cornell LII}).

Turning to the District's Establishment Clause defense in Part IV, the Court rejected the premise that the First Amendment's Clauses are ``warring'' rather than ``complementary'' (\href{https://www.law.cornell.edu/supremecourt/text/21-418}{Cornell LII}). It declared that the Court had ``long ago abandoned \textit{Lemon} and its endorsement test offshoot,'' criticizing them as ``ambitious, abstract, and ahistorical'' tests that invited chaos and functioned as a ``modified heckler's veto'' (\href{https://www.law.cornell.edu/supremecourt/text/21-418}{Cornell LII}). Instead, courts must interpret the Establishment Clause by ``reference to historical practices and understandings'' reflecting the Founding Fathers' understanding (\href{https://www.law.cornell.edu/supremecourt/text/21-418}{Cornell LII}). While acknowledging that government coercion of religious observance is a historical hallmark of establishment, the Court found no evidence of coercion in the record for the three October 2015 prayers: no Bremerton students joined him, the District conceded it had no evidence of direct coercion, and mere visibility or offense ``does not equate to coercion'' (\href{https://www.law.cornell.edu/supremecourt/text/21-418}{Cornell LII}). Distinguishing \textit{Lee v. Weisman}, 505 U.S. 577 (1992), and \textit{Santa Fe Independent School District v. Doe}, 530 U.S. 290 (2000), as involving captive audiences or school-sponsored broadcasts, the Court concluded that ``phantom constitutional violations'' cannot justify actual infringements of First Amendment rights (\href{https://www.law.cornell.edu/supremecourt/text/21-418}{Cornell LII}).

\subsection{Concurring Opinion (Justice Thomas)}
Justice Thomas joined the Court's opinion in full but wrote separately to highlight two unresolved Free Exercise questions: (1) whether public employees' Free Exercise rights differ from those of the general public in light of history and tradition, and (2) what level of scrutiny a government employer must satisfy to restrict an employee's religious expression, noting that the Court has never applied \textit{Pickering} balancing to a Free Exercise claim (\href{https://www.law.cornell.edu/supremecourt/text/21-418}{Cornell LII}).

\subsection{Concurring Opinion (Justice Alito)}
Justice Alito joined the Court's opinion in full on the understanding that it left open what Free Speech standard applies when a public employee engages in purely private expression during a brief lull in official duties, holding only that the District's retaliation failed under any of the discussed standards (\href{https://www.law.cornell.edu/supremecourt/text/21-418}{Cornell LII}).

\subsection{Dissenting Opinion (Justice Sotomayor)}
Justice Sotomayor, joined by Justices Breyer and Kagan, argued that the majority ``misconstrues the facts'' by isolating three October 2015 prayers from Kennedy's seven-year history of leading audible midfield prayers while holding up helmets, his media campaign, and the resulting disruption when spectators rushed the field and knocked over band members (\href{https://www.law.cornell.edu/supremecourt/text/21-418}{Cornell LII}). Under \textit{Engel v. Vitale}, \textit{Abington School District v. Schempp}, \textit{Lee}, and \textit{Santa Fe}, public schools must maintain strict neutrality because students are uniquely susceptible to subtle coercive pressure from coaches who control playing time and college recommendations (\href{https://www.law.cornell.edu/supremecourt/text/21-418}{Cornell LII}). She faulted the majority for overruling \textit{Lemon} in all contexts, discarding the endorsement inquiry, and replacing it with an undefined ``history and tradition'' test that offers ``essentially no guidance for school administrators'' while weakening the Establishment Clause ``backstop'' that protects religious liberty for all (\href{https://www.law.cornell.edu/supremecourt/text/21-418}{Cornell LII}).

\section{Examples: Future Applications}
\begin{enumerate}
    \item \textbf{Same-Side (Protected Private Observance):} A public high school math teacher bows her head and silently prays for one minute at her desk during her scheduled lunch period in the faculty lounge, where students occasionally pass by the open door. Under \textit{Kennedy}, this is protected private religious exercise and speech during personal time, with no coercion or official sponsorship (\href{https://www.law.cornell.edu/supremecourt/text/21-418}{Cornell LII}).
    \item \textbf{Same-Side (Visible Religious Attire):} A public middle school basketball coach wears a yarmulke or hijab on the sidelines during games, while the school permits coaches to wear secular hats or jewelry. Because prohibiting religious attire while allowing comparable secular items is neither neutral nor generally applicable, and visible religious identity is not coercive, \textit{Kennedy} squarely protects the coach (\href{https://www.law.cornell.edu/supremecourt/text/21-418}{Cornell LII}).
    \item \textbf{Opposite-Side (Captive-Audience Official Prayer):} A public high school principal opens mandatory morning homeroom announcements over the school intercom by reciting a Christian prayer and inviting students to bow their heads. Even after \textit{Kennedy}, this violates the Establishment Clause: the principal speaks pursuant to official duties to a captive audience of minors, implicating the coercion and state-sponsored worship barred by \textit{Engel}, \textit{Lee}, and \textit{Santa Fe} (\href{https://www.law.cornell.edu/supremecourt/text/21-418}{Cornell LII}).
    \item \textbf{Opposite-Side (Coach-Led Team Prayer):} A public high school soccer coach gathers the entire team in a mandatory postgame huddle before dismissing them to their parents and leads an audible prayer while instructing players to take a knee. Because the coach exercises official authority over a captive student audience during a required team activity, this crosses the coercion line that \textit{Kennedy} explicitly preserved (\href{https://www.law.cornell.edu/supremecourt/text/21-418}{Cornell LII}).
    \item \textbf{Fence-Sitter (Quiet Prayer Immediately Before Official Huddle):} A public high school track coach kneels silently for 20 seconds at the finish line right after a race, knowing that several varsity runners routinely gather around him to pray before he immediately stands up to deliver the mandatory post-meet critique. This is a close case: defenders would invoke \textit{Kennedy} because the coach neither invites nor requires student participation during a brief lull, while challengers would argue that the immediate transition into an official huddle and recurring student participation create indirect peer and authority coercion under \textit{Lee} and \textit{Santa Fe} (\href{https://www.law.cornell.edu/supremecourt/text/21-418}{Cornell LII}; \href{https://lawreview.uchicago.edu/online-archive/establishment-originalism-kennedy-v-bremerton-school-district}{Univ.\ of Chicago L.\ Rev.\ Online}).
\end{enumerate}

\section{Critique}
\textbf{Originalist and Conservative Perspectives:} Supporters praise \textit{Kennedy} for finally burying the ``ambitious, abstract, and ahistorical'' \textit{Lemon} test and harmonizing the First Amendment's three Clauses rather than treating the Establishment Clause as a trump card over religious expression (\href{https://fedsoc.org/fedsoc-review/a-cord-of-three-strands-how-kennedy-v-bremerton-school-district-changed-free-exercise-establishment-and-free-speech-clause-doctrine}{Federalist Society Review}). Scholars such as Stephanie Taub and Kayla Ann Toney argue that looking to ``historical practices and understandings''---informed by Professor Michael McConnell's six Founding-era hallmarks of establishment cited in Justice Gorsuch's \textit{Shurtleff} concurrence---anchors Establishment Clause doctrine in original public meaning rather than judicial policy preferences (\href{https://fedsoc.org/fedsoc-review/a-cord-of-three-strands-how-kennedy-v-bremerton-school-district-changed-free-exercise-establishment-and-free-speech-clause-doctrine}{Federalist Society Review}). At the same time, some originalist-leaning commentators note that Justice Gorsuch's opinion is surprisingly thin on actual Founding-era analysis and never clearly explains how to analogize eighteenth-century institutional establishments to modern public schools, which did not exist in 1791 (\href{https://lawreview.uchicago.edu/online-archive/establishment-originalism-kennedy-v-bremerton-school-district}{Univ.\ of Chicago L.\ Rev.\ Online}).

\textbf{Progressive and Doctrinal Perspectives:} Critics such as Ira C. Lupu and Robert W. Tuttle condemn \textit{Kennedy} as ``a sledgehammer to the bedrock of nonestablishment,'' arguing that the majority distorted the factual record to portray a publicized, multi-year ritual as a quiet personal moment and effectively hollowed out the post-WWII School Prayer Cases (\href{https://www.gwlr.org/kennedy-v-bremerton-school-district-a-sledgehammer-to-the-bedrock-of-nonestablishment/}{Geo.\ Wash.\ L.\ Rev.\ On the Docket}). They warn that reducing Establishment Clause protection to proof of direct coercion ignores the subtle power dynamics between coaches and adolescents, leaves minority-faith and nonreligious students vulnerable to majoritarian pressure, and gives school administrators virtually no workable standard (\href{https://www.gwlr.org/kennedy-v-bremerton-school-district-a-sledgehammer-to-the-bedrock-of-nonestablishment/}{Geo.\ Wash.\ L.\ Rev.\ On the Docket}; \href{https://www.law.cornell.edu/supremecourt/text/21-418}{Cornell LII}).

\textbf{Underlying Competing Values:} The case exposes a deep structural tension between two visions of religious liberty in a pluralistic democracy: (1) an \textit{accommodationist} vision prioritizing individual religious expression and equal treatment of religious and secular speech in the public square, versus (2) a \textit{separationist} vision prioritizing institutional neutrality in compulsory public schools to shield impressionable students from state-backed religious conformity (\href{https://constitutioncenter.org/the-constitution/supreme-court-case-library/kennedy-v-bremerton-school-district}{National Constitution Center}).

\section{Key Quotations}
\begin{quote}
``The Constitution and the best of our traditions counsel mutual respect and tolerance, not censorship and suppression, for religious and nonreligious views alike.''
\end{quote}
--- Justice Neil Gorsuch, Opinion of the Court (\href{https://www.law.cornell.edu/supremecourt/text/21-418}{Cornell LII}; \href{https://www.courtlistener.com/opinion/6619004/kennedy-v-bremerton-school-dist/}{CourtListener}).

\begin{quote}
``In place of \textit{Lemon} and the endorsement test, this Court has instructed that the Establishment Clause must be interpreted by `` `reference to historical practices and understandings.' '' ''
\end{quote}
--- Justice Neil Gorsuch, Opinion of the Court (\href{https://www.law.cornell.edu/supremecourt/text/21-418}{Cornell LII}).

\begin{quote}
``A natural reading of that sentence would seem to suggest the Clauses have ``complementary'' purposes, not warring ones where one Clause is always sure to prevail over the others.''
\end{quote}
--- Justice Neil Gorsuch, Opinion of the Court (\href{https://www.law.cornell.edu/supremecourt/text/21-418}{Cornell LII}).

\begin{quote}
``And in no world may a government entity's concerns about phantom constitutional violations justify actual violations of an individual's First Amendment rights.''
\end{quote}
--- Justice Neil Gorsuch, Opinion of the Court (\href{https://www.law.cornell.edu/supremecourt/text/21-418}{Cornell LII}).

\begin{quote}
``Today's decision is particularly misguided because it elevates the religious rights of a school official, who voluntarily accepted public employment and the limits that public employment entails, over those of his students, who are required to attend school and who this Court has long recognized are particularly vulnerable and deserving of protection.''
\end{quote}
--- Justice Sonia Sotomayor, Dissenting Opinion (\href{https://www.law.cornell.edu/supremecourt/text/21-418}{Cornell LII}).

\end{document}
